
\documentclass[aspectratio=169]{beamer}

\mode<presentation>
\usetheme{Boadilla}
\definecolor{NTHU1}{RGB}{127,16,132}
\definecolor{NTHU2}{RGB}{228,210,231}
\definecolor{blue}{RGB}{30,90,205}
\definecolor{red}{RGB}{213,94,0}
\definecolor{green}{RGB}{0,128,0}
\definecolor{charcoalgray}{RGB}{108,104,104}
\setbeamercolor{title}{fg=NTHU1}
\setbeamercolor{frametitle}{fg=NTHU1}
\setbeamercolor{block title}{bg=NTHU1, fg=white}
\setbeamercolor{block body}{bg=NTHU2}
\setbeamercolor{structure}{fg=NTHU1}
\setbeamercolor{item projected}{fg=white}
\setbeamercolor{item}{fg=NTHU1}
\setbeamercolor{subitem}{fg=NTHU1}
\setbeamercolor{section in toc}{fg=NTHU1}
\setbeamercolor{description item}{fg=NTHU1}
\setbeamercolor{caption name}{fg=NTHU1}
\setbeamercolor{button}{bg=NTHU1, fg=white}
\setbeamercolor{caption name}{fg=NTHU1}
\usepackage{graphics}
\usepackage{tikz}
\usepackage{amsmath}
\usepackage{bbm}
\usetikzlibrary{decorations.pathreplacing}
\usepackage{geometry}
\usepackage{booktabs}
\usepackage{bookmark}
\usepackage{multirow, makecell}
\usepackage{float}
\usepackage{fancyvrb}
\usepackage{caption}
\captionsetup[figure]{singlelinecheck = false, format= hang, justification=centering}
\captionsetup[subfigure]{singlelinecheck = false, format= hang, justification=raggedright}
\usepackage{subcaption}
\usepackage{adjustbox}
\usepackage{hyperref}
\usepackage{threeparttable}
\usepackage{appendixnumberbeamer} 
\usepackage[T1]{fontenc}
\usepackage[default]{lato} 
\usepackage{smartdiagram}
\smartdiagramset{
  back arrow disabled = true,
  module minimum width=1cm,
  module x sep = 3,
  uniform arrow color=true,
}
\usepackage{xcolor}
\usepackage{colortbl}
  \definecolor{light-gray}{gray}{0.99}

\newenvironment{wideitemize}{\itemize\addtolength{\itemsep}{10pt}}{\enditemize}
\newenvironment{wideenumerate}{\enumerate\addtolength{\itemsep}{10pt}}{\endenumerate}
\newenvironment{widedescription}{\description\addtolength{\itemsep}{10pt}}{\enddescription}
\hypersetup{
colorlinks=true,
linkcolor=NTHU1,
filecolor=green, 
urlcolor=blue,
}
\beamertemplatenavigationsymbolsempty
\setbeamercolor{author in head/foot}{bg=white, fg=NTHU1}
\setbeamercolor{title in head/foot}{bg=white, fg=NTHU1}
\setbeamercolor{date in head/foot}{bg=white, fg=NTHU1}
\setbeamercolor{section in head/foot}{bg=white, fg=NTHU1}
\setbeamercolor{headline}{bg=white}
\setbeamertemplate{footline}{
    \leavevmode%
    \hbox{%
        \begin{beamercolorbox}[wd=.333333\paperwidth,ht=2.25ex,dp=1ex,center]{date in head/foot}%
            \usebeamerfont{date in head/foot}%\insertdate
        \end{beamercolorbox}%
        \begin{beamercolorbox}[wd=.333333\paperwidth,ht=2.25ex,dp=1ex,center]{title in head/foot}%
            \usebeamerfont{title in head/foot}\insertshorttitle
        \end{beamercolorbox}%
        \begin{beamercolorbox}[wd=.333333\paperwidth,ht=2.25ex,dp=1ex,right]{date in head/foot}%
            \usebeamerfont{date in head/foot}\insertframenumber{}\hspace*{2em}
        \end{beamercolorbox}}%
        \vskip0pt%
    }
    %% May need to script the introduction!!!!!!
%\setbeamercolor{page number in head/foot}{fg=NTHU1}
\setbeamertemplate{section in toc}[sections numbered]
\setbeamertemplate{subsection in toc}{\leavevmode\leftskip=3em\rlap{\hskip-1.75em\inserttocsectionnumber.\inserttocsubsectionnumber}
\inserttocsubsection\par}
\setbeamerfont{subsection in toc}{size=\footnotesize}
%\setbeamertemplate{headline}{%
  %\begin{beamercolorbox}[ht=5.5ex]{section in head/foot}
    %\vskip2pt\insertnavigation{0.33\paperwidth}\vskip2pt
  %\end{beamercolorbox}%
%}
\newenvironment{transitionframe}{\setbeamercolor{background canvas}{bg=white}\setbeamertemplate{footline}{} \begin{frame}[noframenumbering]}{\end{frame}}

\makeatletter
\let\@@magyar@captionfix\relax
\makeatother

\title[]{Conflict and Peace} % Change this regularly
\subtitle[]{Week 7: Impact of crime and conflict on economic and institutional development}
\author[Lee]{Seung-hun Lee }
\institute[]{Taipei School of Economics\\ National Tsing Hua University}

\date[]{October 17th, 2025}

\begin{document}
\begin{transitionframe}
\titlepage
\end{transitionframe}

%%% Color slides for section headers: Use for colloquium version (The ones bewteen \iffals and \fi)

\section{Overview of the topic}
\begin{frame}
\frametitle{How do conflicts affect economic and institutional development?} 
Conflicts impose negative impact both in short and long-run
\begin{itemize}\setlength\itemsep{3pt}
\item Short: Loss of human life, destruction of capital, increased uncertainty
\item Long: Loss of tax base $\to$ Difficult to finance reconstruction, reliance on debt $\Uparrow$
\item It can also spill over into other regions (depress trade, increased social strains)
\end{itemize}\medskip
Conflicts undermine the foundation for institutions that support growth and development
\begin{itemize}\setlength\itemsep{3pt}
\item Weakened governance: Perpetrators gain influence over decisionmakers
\item Weakened accountability: General public gets less priority
\item Leads to more corrupt and incapable state
\end{itemize}\medskip
These factors put many conflict-affected societies into ``conflict trap"
\begin{itemize}\setlength\itemsep{3pt}
\item Conflict $\to$ Economy and institution$\Downarrow$ $\to$ Conflict $\to$Economy and institution$\Downarrow$ $\to$.....
\end{itemize}\medskip
\end{frame}

\begin{frame}
\frametitle{Long-run effects of armed conflicts} 
\begin{figure}
\begin{subfigure}{0.475\textwidth}
  \includegraphics[keepaspectratio,width=\textwidth]{fig1.png}
  \caption{Actual GDP vs GDP with no conflict}
\end{subfigure}
\begin{subfigure}{0.45\textwidth}
  \includegraphics[keepaspectratio,width=\textwidth]{fig2.png}
  \caption{Accrual of debt post-conflict}
\end{subfigure}
\end{figure}
\end{frame}


\begin{frame}
\frametitle{Can scars of conflict and crime be undone?} 
Wartime trauma has serious long-run development consequences
\begin{itemize}\setlength\itemsep{3pt}
\item Stunted economic development and recurrence of war if trauma is not treated
\item But also difficult to escape: Many reconstruction programs do face shortfalls
\begin{itemize}\setlength\itemsep{3pt}
\item Many reconstruction fail to foster trust across fractionalized groups and outsiders
\item Reintegration programs sometimes fail to generate interest among participants
\end{itemize}

\item So what are the required approaches
\begin{itemize}\setlength\itemsep{3pt}
\item Reshape incentives: Make it more costly to engage in conflicts
\item Winning hearts and minds: Trust across warring groups need to be rebuilt
\item For intervention led by third party: Trust in intervening entities too
\end{itemize}
\item Some of these efforts do succeed in rebuilding capital and trust
\end{itemize}
\end{frame}

\begin{frame}
\frametitle{Some countries do find some hope} 
\begin{figure}
  \includegraphics[keepaspectratio,width=0.7\textwidth]{fig3.png}
\end{figure}
\end{frame}

\begin{frame}
\frametitle{Question to explore today} 

\begin{itemize}\setlength\itemsep{3pt}
\item How does conflict hurt aggregate economic and institutional performance?
\begin{itemize}\setlength\itemsep{3pt}
\item Destroys foundations needed for economic activities and growth
\item Distorts resource allocation to favor attackers, not general public
\item Captures state apparatus (politicians, institutions) $\to$ lead to above consequences
\end{itemize}
\item How do you design effective counter-conflict programs?
\begin{itemize}\setlength\itemsep{3pt}
\item Just addressing perpetrators may not suffice (e.g Mexico)
\item Need to reshape incentives and perceptions
\end{itemize}
\end{itemize} \medskip
We investigate how conflicts affect economic and institutional development \\ \medskip
We then probe into how conflicts can be undone
\end{frame}

\begin{transitionframe}
\frametitle{Roadmap for today}
\tableofcontents
\end{transitionframe}

\section{Economic impact of conflicts}
\subsection{Pinotti (2015): The Economic Costs of Organised Crime: Evidence from Southern Italy}
\begin{transitionframe}
\begin{center}
  \begin{huge}
    \textcolor{NTHU1}{%
Pinotti (2015):  \\ The Economic Costs of Organised Crime: Evidence from Southern Italy
    }
  \end{huge}
\end{center}
\end{transitionframe}

\begin{frame}
\frametitle{Q: Presence of organized crime $\Rightarrow$ economic development?} 
Organized crime is perceived as barrier to economic development
\begin{itemize}\setlength\itemsep{3pt}
\item Not just in developing countries, but also more developed countries too 
\item Among developed countries, Italy stands out
\begin{itemize}\setlength\itemsep{3pt}
\item 5 regions with highest organized crime are also the poorest
\end{itemize}
\item Challenge: Reverse causality among economic development and organized crime
\begin{itemize}\setlength\itemsep{3pt}
\item We want to confirm organized crime $\Uparrow$ $\Rightarrow$ economic development $\Downarrow$
\item It could be organized crime $\Uparrow$ $\Leftarrow$ economic development $\Downarrow$
\end{itemize}
\end{itemize}\medskip
This paper addresses this with unexpected changes in criminal presence in select regions
\begin{itemize}\setlength\itemsep{3pt}
\item Setting: Apulia and Basilicata (normal levels of criminal presence $\to$ heavy)
\item This trend arose due to changes in entry point of smuggled tobacco in Italy
\item These changes are unrelated to socio-economic environments
\item Generate counterfactual using synthetic control methods
\end{itemize}
\end{frame}


\begin{frame}
\frametitle{History of organized crime in Italy} 
Description of their activities
\begin{itemize}\setlength\itemsep{3pt}
\item  Supply illicit goods, extortion and threats, private protection where state is absent
\item Arose from the vacuum property rights protection following 1861 Italian Unification
\item Italian judicial system began to address organized crime actively since 1980s (!)
\begin{itemize}
  \item Before: Applied rules that treats organized crime equally with small \# of individuals
\end{itemize}
\end{itemize}\medskip
Economic growth and crime in Southern Italy
\begin{itemize}\setlength\itemsep{3pt}
  \item Generally, convergence between north and south stopped around 70s
  \item Apulia and Basilicata: High growth in 60s-70s $\to$ lowest growth rates in Italy
  \item The down-growth period coincides with increased organized criminal presence
 \begin{itemize} 
  \item Tobacco smuggling $\Uparrow$ + Relief money after major earthquake + Criminals \textit{in confino}
\end{itemize}
\end{itemize}
\end{frame}

\begin{frame}
\frametitle{Data and empirical strategy} 
Measuring Crime
\begin{itemize}\setlength\itemsep{3pt}
\item Data from Parliamentary Antimafia Commission and Italian Statistical institute
\item Regional GDP per capita from 1951-2005
\item Value added by industry and population density accounted for
\end{itemize}\medskip
Synthetic control to create a counterfactual of the treated areas
\begin{itemize}\setlength\itemsep{3pt}
  \item Idea: Actual outcome in treated area vs weighted average of control group
  \item Weights are chosen to match pre-treatment characteristics of treated areas
  \item Includes other regions except Sicily, Campania, and Calabria (pervasive OCG presence) 
  \item Therefore, this method generates counterfactual out of the basket of control group
\end{itemize}
\end{frame}

\begin{frame}
\frametitle{Trend in growth and organized criminal presence} 
\begin{figure}
  \begin{subfigure}{0.35\textwidth}
  \includegraphics[keepaspectratio, width=\textwidth]{p2015_f1.jpeg}
  \caption{OCG distribution}
  \end{subfigure}
  \begin{subfigure}{0.6\textwidth}
  \includegraphics[keepaspectratio, width=\textwidth]{p2015_f2.jpeg}
  \caption{Growth rates across time and regions}
  \end{subfigure}
\end{figure}
\end{frame}

\begin{frame}
\frametitle{Criminmal presence negatively affects growth} 
\begin{figure}
  \begin{subfigure}{0.475\textwidth}
  \includegraphics[keepaspectratio, width=\textwidth]{p2015_t1.jpeg}
  \caption{Synthetic control similar to treatment}
  \end{subfigure}
  \hfill
  \begin{subfigure}{0.475\textwidth}
  \includegraphics[keepaspectratio, width=\textwidth]{p2015_f3.jpeg}
  \caption{Actual GDP $<$ Synthetic GDP growth}
  \end{subfigure}
\end{figure}
\end{frame}

\begin{frame}
\frametitle{Results largely driven by decrease in economic activity} 
\begin{figure}
  \includegraphics[keepaspectratio, width=0.65\textwidth]{p2015_f4.jpeg}
\end{figure}
\end{frame}

\begin{frame}
\frametitle{Discussion and takeaways} 
Other mechanisms tested
\begin{itemize}\setlength\itemsep{3pt}
\item No shift in economic activities from formal to informal sectors
\item Measured with electricity usage (even informal sectors use them)
\item Electricity decreases: Suggests that the economic activities themselves decrease
\end{itemize}\medskip
Takeaways
\begin{itemize}\setlength\itemsep{3pt}
  \item Presence of organized crime $\to$ Economic development $\Downarrow$
  \item Driven by decrease in economic activities (investment, electricity usage)
  \item Macroeconomic outcomes: Does not capture effects on individual economic activities
  \begin{itemize}\setlength\itemsep{3pt}
  \item Need to observe individuals as unit of analysis (last week!)
  \end{itemize}
  \item Also, there could be evidence that informal sectors grow due to organized crime
\end{itemize}
\end{frame}

\subsection{Fenizia, Saggio (2024): Organized Crime and Economic Growth: Evidence from Municipalities Infiltrated by the Mafia}
\begin{transitionframe}
\begin{center}
  \begin{huge}
    \textcolor{NTHU1}{%
Fenizia, Saggio (2024): \\ Organized Crime and Economic Growth: Evidence from Municipalities Infiltrated by the Mafia
    }
  \end{huge}
\end{center}
\end{transitionframe}

\begin{frame}
\frametitle{Q: Break up institutions infiltrated by crime $\Rightarrow$ development?} 
How do you undo effects of organized crime?
\begin{itemize}\setlength\itemsep{3pt}
\item Illegal activities, extortions, hinders competition, stunt development
\item What are the effects of efforts to regain control from organized crime?
\begin{itemize}
\item Ultimately relates to regaining legitimacy in areas with long history of organized crime
\end{itemize}
\end{itemize} \medskip
This paper analyzes the long-run economic effects of actions to regain control
\begin{itemize}\setlength\itemsep{3pt}
\item Italy: Dismiss city councils infiltrated by the Mafia (CCD)
\begin{itemize}\setlength\itemsep{3pt}
\item Central government dismiss mayor and city council and appoints commissioners
\item Commissioners administer with full legislative and executive authorities until next election 
\end{itemize}
\item Compare areas with CCDs to observationally similar municipalities w/o CCDs
\item Leverage changes in economic activities: Employment, Firms, Real estate price etc
\end{itemize}
\end{frame}

\begin{frame}
\frametitle{Counter-organized crime activities in Italy} 
Dismissal of city councils (CCDs)
\begin{itemize}\setlength\itemsep{3pt}
\item Introduced in 1991 in response to rise in Mafia influence in 1980s
\item Any sign of Mafia influence detected $\to$ appoint commissioner to replace city council
\begin{itemize}\setlength\itemsep{3pt}
\item Usually instigated by unrelated police investigations 
\item Lower evidentiary bar compared to legal prosecution
\end{itemize}
\item The commissioners are experienced career civil servants
\item Legislative and executive authority until next local elections (roughly 24-36 months)
\end{itemize} \medskip
 What do these commissioners do?
\begin{itemize}\setlength\itemsep{3pt}
\item Freeze all investment in new projects (review financial situation)
\item Revoke procurement contracts, permits, licenses with connection to organized crime
\item Replace municipal personnel and mandate training 
\item Attempt to gain trust and support of local communities (free job training etc)
\end{itemize}
\end{frame}

\begin{frame}
\frametitle{Data and empirical strategy} 
Sources of data
\begin{itemize}\setlength\itemsep{3pt}
\item Employer-Employee matched social security data: Employment, wage, residence etc
\item Real estate prices from Italian Treasury
\item Politicians data from Ministry of interior (educational attainment and age)
\item Councils: Ministry of interior (expenditures, revenues) \& ANAC (contracts)
\end{itemize} \medskip
Leverage DID matching treated area to observationally similar control group (Matched DID)
\[
Y_{mt}=\alpha_m + \lambda_{r(m),t}+\sum_{k=a}^b \tilde{\theta}^k I[t=t_m^*+k]+\sum_{k=a(\neq-1)}^b \theta^k I[t=t_m^*+k]\times \text{CCD}_m + u_{mt}
\]\vspace{-10pt}
\begin{itemize}\setlength\itemsep{3pt}
\item Each treated municipalities are matched to close counterparts without CCD
\item $t_m^*$: Year of CCD for municipality $m$, $I[t=t_m^*+k]$ is the event-time dummy
\item $\tilde{\theta}^k$: Effect of event time dummies for control group (similar to time FE)
  \item $\theta^k$ is the effect of CCD $k$ years later on the treated municipalities
\end{itemize}
\end{frame}

\begin{frame}
\frametitle{CCDs over time and place} 
\begin{figure}
\begin{subfigure}{0.375\textwidth}
\includegraphics[keepaspectratio,width=\textwidth]{fs2024_f1_1.jpeg}
\end{subfigure}
\begin{subfigure}{0.575\textwidth}
\includegraphics[keepaspectratio,width=\textwidth]{fs2024_f1_2.jpeg}
\end{subfigure}
\end{figure}
\end{frame}

\begin{frame}
\frametitle{Rising employment and stock of firms} 
\begin{figure}
\includegraphics[keepaspectratio,width=\textwidth]{fs2024_f2.jpeg}
\end{figure}
\begin{itemize}
  \item Industrial real estate prices also rises (not on this slide)
\end{itemize}
\end{frame}

\begin{frame}
\frametitle{Driven by weakened mafia connections} 
\begin{figure}
\begin{subfigure}{0.525\textwidth}
\includegraphics[keepaspectratio,width=\textwidth]{fs2024_f3.jpeg}
\caption{New face to local politics}
\end{subfigure}
\hspace{20pt}
\begin{subfigure}{0.275\textwidth}
\includegraphics[keepaspectratio,width=\textwidth]{fs2024_f4.jpeg}
\caption{Old corrupt politicians out}
\end{subfigure}
\end{figure}
\end{frame}

\begin{frame}
\frametitle{Takeaways and discussions} 
Key takeaways
\begin{itemize}\setlength\itemsep{3pt}
\item Dismissing Mafia influence allows formal governments to reassert legitimacy 
\item Long term effects $>$ Short-run effects
\item New faces: Suggests that Mafia presence is actually diminishing here
\item Effects of directly targeting local institutions $>$ just illicit activities (see Mexico)
\end{itemize} \medskip
Discussions for further studies
\begin{itemize}\setlength\itemsep{3pt}
  \item Institutional shake-up is required to counter long-standing crime/conflict?
\item Why did the Mafia not fight back?
\begin{itemize}\setlength\itemsep{3pt}
  \item They may have changed their operation strategy (less confrontation with the state)
  \item Would this work in places where organized crime actively confronts the state?
\end{itemize}
\end{itemize}
\end{frame}

\subsection{Korovkin, Marakin (2023): Conflict and Intergroup Trade: Evidence from the 2014 Russia-Ukraine Crisis [student presentation]}
\begin{transitionframe}
\begin{center}
  \begin{huge}
    \textcolor{NTHU1}{%
Korovkin, Marakin (2023): \\ Conflict and Intergroup Trade: Evidence from the 2014 Russia-Ukraine Crisis [student presentation]
    }
  \end{huge}
\end{center}
\end{transitionframe}



\section{Institutional impact of conflicts}
\subsection{Alesina, Piccolo, Pinotti (2019): Organized Crime, Violence, and Politicians}
\begin{transitionframe}
\begin{center}
  \begin{huge}
    \textcolor{NTHU1}{%
Alesina, Piccolo, Pinotti (2019): \\ Organized Crime, Violence, and Politicians
    }
  \end{huge}
\end{center}
\end{transitionframe}

\begin{frame}
\frametitle{Q: How do organized crime influence institutional quality?} 
One of the worst aspects of organized crime: State capture
\begin{itemize}\setlength\itemsep{3pt}
\item Distort allocation of resources in favor of the captors (criminal organizations)
\item Unjustly relax rules for their favor
\item Why and how do criminal organizations take over politicians and elections?
\end{itemize}\medskip
This paper models this phenomena and identifies empirical patterns 
\begin{itemize}\setlength\itemsep{3pt}
\item Violence: Signal strength and will to use violence 
\item Candidates scared off and influences behavior of elected politicians
\item Leverages presence of Sicilian Mafia to test theoretical predictions
\item Compare regions with and without significant presence of organized crime
\end{itemize}
\end{frame}

\begin{frame}
\frametitle{Organized crime and politics in Italy} 
Key features of organized criminal group (Italy, but also others)
\begin{itemize}\setlength\itemsep{3pt}
\item Command armed men capable of exerting intimidation and influence official economy
\item Sicilian Mafia colluded with certain parties
\begin{itemize}
\item Christian Democrats (CD) accepted Mafia support to gain advantage against opponents
\end{itemize}
\item Targeted political activists and politicians supporting left-wing parties
\item Target national elections: National officials determine scope of criminal group activities
\end{itemize}\medskip
Theoretical predictions that explain criminal organization behaviors
\begin{itemize}\setlength\itemsep{3pt}
\item More likely to be violent during elections, and in swing districts (majoritarian system)
\item Violence results in lower votes for honest party, higher votes for captured parties
\item Anti-organized crime effort $\to$ retaliation $\to$ lower future anti-crime efforts
\item These efforts from honest party decrease with rising pre-election crimes
\end{itemize}

\end{frame}

\begin{frame}
\frametitle{Data and empirical strategy} 
Data sources
\begin{itemize}\setlength\itemsep{3pt}
\item Criminal group presence: Various NGOs (e.g Libera)
\item Electoral results from Italian Interior Ministry
\item Politician behavior: Transcripts of speeches from 1948-2008 in national Parliament
\begin{itemize}
  \item Track mentions of Mafia-related terms (almost all of them in negative light)
\end{itemize}
\item 
\end{itemize} \medskip
Empirical strategy leverages differences in Mafia vs non-Mafia regions
\begin{itemize}
  \item Primarily Difference-in-differences
  \item But also does wide range of heterogeneity analysis (Triple differences)
\[
  \begin{aligned}
  Y_{it} &= \alpha+\beta_1\text{Treat1}_i + \beta_2\text{Treat2}_i + \beta_3\text{Post}_t  \\
         & + \beta_4\text{Treat1}_i \times \text{Treat2}_i + \beta_5 \text{Treat1}_i \times \text{Post}_t + \beta_6\text{Treat2}_i \times \text{Post}_t  \\
         & +\beta_7\text{Treat1}_i\times \text{Treat2}_i \times \text{Post}_t +e_{it}
  \end{aligned}
  \]

\end{itemize}
\end{frame}

\begin{frame}
\frametitle{More attacks, particularly in swing states}
\begin{figure}
  \includegraphics[keepaspectratio, width=0.67\textwidth]{app2019_t1.png}
\end{figure}
\end{frame}

\begin{frame}
\frametitle{Uncooperative party vote share falls}
\begin{figure}
  \includegraphics[keepaspectratio, width=0.65\textwidth]{app2019_t2.png}
\end{figure}
\end{frame}

\begin{frame}
\frametitle{Political attacks discourage politicians (esp left-wing)}
\begin{figure}
  \includegraphics[keepaspectratio, width=0.65\textwidth]{app2019_t3.png}
\end{figure}
\end{frame}

\begin{frame}
\frametitle{Discussion and Takeaways} 
Criminal organizations and conflicts capture politicians and the states
\begin{itemize}\setlength\itemsep{3pt}
\item Strategically use elections: Timing and place
\item Aim to favor the cooperating politicians by hurting the uncooperative
\item Does affect behaviors of elected politicians: Reduce anti-conflict efforts
\item They especially affect efforts of the honest politicians (capture ones are unaffected)
\end{itemize}\medskip
Discussions
\begin{itemize}\setlength\itemsep{3pt}
\item Generalizable? Likely (we have already seen this with insurgents in Afghanistan)
\item Long run effects of such capture? (Distortionary allocations?)
\begin{itemize}\setlength\itemsep{3pt}
  \item How are public expenditures and procurement affected?
  \item These could skew towards favoring captors over general public
  \item We look extensively at weak states and state capture late in the semester
\end{itemize}
\end{itemize}
\end{frame}


\subsection{Brückner, Ciccone (2011): Rain and the Democratic Window of Opportunity }
\begin{transitionframe}
\begin{center}
  \begin{huge}
    \textcolor{NTHU1}{%
Brückner, Ciccone (2011): \\ Rain and the Democratic Window of Opportunity 
    }
  \end{huge}
\end{center}
\end{transitionframe}

\begin{frame}
\frametitle{Q: Can economic shocks induce democratic transitions?} 
What causes movements behind democratic change?
\begin{itemize}\setlength\itemsep{3pt}
\item Many argue that economic recessions lead to democratic improvements
\begin{itemize}\setlength\itemsep{3pt}
\item Transitory shocks $\to$ window of opportunity to contest power
\item Autocratic regime may have to make concessions
\end{itemize}
\item Challenges: Democratic transition affect economics (reverse causality)
\end{itemize}\medskip
This paper isolates exogenous link between recessions and democratic transitions
\begin{itemize}\setlength\itemsep{3pt}
\item Uses negative rainfall shocks (i.e. droughts) that affect economic performance
\item Finds that even temporary shocks can lead to concessions from autocratic regimes
\item Rainfall $\Downarrow$ $\to$ GDP \& opportunity cost of contesting power $\Downarrow$ $\to$ Democratic transitions
\end{itemize}
\end{frame}

\begin{frame}
\frametitle{Dataset and definitions} 
Sources of main measures and definitions
\begin{itemize}\setlength\itemsep{3pt}
\item Polity IV: Polity2 score that defines level of democracy (1984-2004)
\begin{itemize}\setlength\itemsep{3pt}
\item Executive constraint + competitive executive recruitment + political participation
\item Range: -10 to 10 (0-: Autocracy, 1-6: Partial democracy, 7-10: Full democracy )
\item Dataset is available as Polity V (updated, \hyperlink{https://www.systemicpeace.org/polityproject.html}{check here}) 
\end{itemize}
\item Democratic transitions at year $t$: Autocracy to democracy between year $t$ and $t+1$
\begin{itemize}
\item Vice versa for autocratic transitions at year $t$
\end{itemize}
\item Democratic step at year $t$: Becomes partial/full democracy between year $t$ and $t+1$
\item Rainfall estimates: NASA Global Precipitation Climatology Project 
\end{itemize}
\end{frame}

\begin{frame}
\frametitle{Democratic transitions in sample countries} 
\begin{figure}
\includegraphics[keepaspectratio,width=0.8\textwidth]{bc2011_f1.png}
\end{figure}
\end{frame}

\begin{frame}
\frametitle{Empirical strategy} 
How to isolate exogenous changes to economic conditions
\begin{itemize}\setlength\itemsep{3pt}
\item Rainfall shock: Affect income through agricultural productivity
\item When income is accounted for, rainfall $\nRightarrow $ democratic transition on its own
\item These constitute relevancy and exogeneity conditions for instrumental variables
\end{itemize}\medskip
Regression specification
\[
\Delta D_{ct} = \alpha_c + \phi_c \times t + \delta _t + \beta_0 \log{y_{ct}}+\beta_1 \widehat{\log{y_{ct-1}}}+\epsilon_{ct}
\]
with the first stage 
\[
\log{y_{ct-1}} = a_c + f_c \times (t-1) + d_{t-1} + b \log{\text{rainfall}_{ct-1}}+e_{ct-1}
\]
\begin{itemize}
\item Why $t-1$: Time it takes from shock to action 
\end{itemize}
\end{frame}

\begin{frame}
\frametitle{Democratic transitions following negative rainfall shocks} 
\begin{figure}
\includegraphics[keepaspectratio,width=0.7\textwidth]{bc2011_t1.png}
\end{figure}
\end{frame}

\begin{frame}
\frametitle{Democratic transitions following negative rainfall shocks} 
\begin{figure}
\includegraphics[keepaspectratio,width=0.7\textwidth]{bc2011_t2.png}
\end{figure}
\end{frame}

\begin{frame}
\frametitle{Takeaways and discussions} 
Economic shocks induced by (lack of) rainfall $\to$ democratic transition
\begin{itemize}\setlength\itemsep{3pt}
  \item Suggests that transitory shocks open up opportunities for change (protests)
  \item These apply to countries reliant on agriculture
  \begin{itemize}
    \item In countries with low agriculture, these effects do not appear
  \end{itemize}
\end{itemize} \medskip
Discussions 
\begin{itemize}\setlength\itemsep{3pt}
  \item Connects to how economic conditions shape incentives to contest power
  \item Also show that these changes can lead to structural and institutional changes
  \item Can this be extended to other exogenous sources to economic shocks?
  \begin{itemize}\setlength\itemsep{3pt}
  \item Sudden disruption of trade opportunities
  \item Resource shocks from foreign countries (esp. for countries importing resources)
  \end{itemize}
  \end{itemize}
\end{frame}



\subsection{Blouin, Murkand (2019):  Erasing Ethnicity? Propaganda, Nation Building, and Identity in Rwanda}
\begin{transitionframe}
\begin{center}
  \begin{huge}
    \textcolor{NTHU1}{%
Blouin, Murkand (2019): \\ Erasing Ethnicity? Propaganda, Nation Building, and Identity in Rwanda
    }
  \end{huge}
\end{center}
\end{transitionframe}

\begin{frame}
\frametitle{Q: Can you undo ethnic divisions that cause conflicts?} 
How do you bridge divides in conflict-ridden society?
\begin{itemize}\setlength\itemsep{3pt}
\item Conflicts, institutional fragility, low growth more likely in countries with ethnic divide
\item Need to increase ethnic trust, cooperation \& reduce salience of ethnics
\begin{itemize}
\item Conversely, nation-building that addresses these points may be effective
\end{itemize}
\end{itemize}\medskip
This paper examines how propaganda can be used in nation-building
\begin{itemize}\setlength\itemsep{3pt}
  \item Propaganda: Govt-owned Radio Rwanda messages on reconciliation
  \begin{itemize}
  \item Contrast to hate message spread by RTLM (Yanagizawa-Drott 2014)
  \end{itemize}
\item Contention as to whether inter-ethnic reconciliation brought real progress 
\begin{quote}
  Mr. Kagame has created a nation that is orderly but repressive ... Against this backdrop, it is difficult to gauge sentiment about the effectiveness of reconciliation efforts \\ (New York Times, 2017)
\end{quote}
\item Leverage geographical traits that determines the broadcast reception
\end{itemize}\medskip
\end{frame}

\begin{frame}
\frametitle{Ethnicity and post-civil war nation-building in Rwanda} 
Post-war Rwanda
\begin{itemize}\setlength\itemsep{3pt}
\item  Three ethnic groups: Hutu, Tutsis (both 97\%), and Twa 
\item Tutsi superiority over Hutus during colonial period $\to$ inter-ethnic resentment $\Uparrow$
\item 1994 Genocide: 70\% of Tutsis slaughtered 
\item Post-conflict: RPF (Tutsi-controlled) asserted control $\to$ Paul Kagame took presidency
\end{itemize}\medskip
Reconciliation under Kagame Administration
\begin{itemize}\setlength\itemsep{3pt}
\item Promote new Rwandan identity: Media, textbooks, and enforced social interaction
\item \textit{Umuganda}: Get individuals from all ethnicities for a collective project on public goods
\item Messages were spread using radio (Radio Rwanda)
\end{itemize}\medskip
\end{frame}

\begin{frame}
\frametitle{Data and empirical strategy} 
Data sources (lab-in-the-field)
\begin{itemize}\setlength\itemsep{3pt}
\item Salience of Identity Test (SIT) score from participants
\item Games to identify inter-ethnicity cooperation: Partner selection game and trust game
\item Individual data: Demographics (not ethnicity - illegal to ask), socioeconomic status
\item Radio signal data from various geospatial sources
\end{itemize}\medskip
Empirical strategy uses randomization and variation across areas
\begin{itemize}\setlength\itemsep{3pt}
\item Leverages differences in radio reception across 27 transmission towers
\[
Y_{ivd}=\alpha_0 + \phi_d + \beta\text{Radio}_{vd}+\Gamma X_{ivd}+e_{ivd}
\]
\item $\text{Radio}_{vd}$: 1 if signal strength $> 45 dB$ at village $v$ in district $d$
\item Outcome at individual level
\end{itemize}\medskip
\end{frame}

\begin{frame}
\frametitle{Variation in radio strength across locations} 
\begin{figure}
\includegraphics[keepaspectratio,width=0.45\textwidth]{bm2019_f1.jpeg}
\end{figure}
\end{frame}

\begin{frame}
\frametitle{Lower salience following exposure to radio} 
\begin{figure}
\includegraphics[keepaspectratio,width=0.8\textwidth]{bm2019_t1.jpeg}
\end{figure}
\end{frame}

\begin{frame}
\frametitle{More inter-ethnic trust and cooperation} 
\begin{figure}
\begin{subfigure}{0.4\textwidth}
\includegraphics[keepaspectratio,width=\textwidth]{bm2019_t2.jpeg}
\caption{Partnership game}
\end{subfigure}
\hspace{20pt}
\begin{subfigure}{0.5\textwidth}
\includegraphics[keepaspectratio,width=\textwidth]{bm2019_t3.jpeg}
\caption{Trust game}
\end{subfigure}
\end{figure}
\end{frame}

\begin{frame}
\frametitle{Takeaways and discussions} 
Key takeaways
\begin{itemize}\setlength\itemsep{3pt}
\item Reconciliation program made some significant progress
\begin{itemize}
\item Salience of ethnicity $\downarrow$, Trust across ethnic groups $\uparrow$
\end{itemize}
\item Nation-building that address inter-ethnic relationship may be productive
\end{itemize} \medskip
Discussions for further studies
\begin{itemize}\setlength\itemsep{3pt}
  \item Is this effect long-lasting? Unsure
   \begin{itemize}
\item Would need a follow-up investigation 
\end{itemize}
\item What happens if the focus of propaganda shifts? Unsure
   \begin{itemize}\setlength\itemsep{3pt}
\item Current reconciliation is the personal agenda of President Kagame
\item What happens if priority shifts?
\end{itemize}
\item Generalizable outside of autocratic (i.e. lack of media freedom) settings? Unsure
\end{itemize}
\end{frame}


\subsection{Berrebi, Klor (2008): Are Voters Sensitive to Terrorism? Direct Evidence from Israeli Electorate [student presentation]}
\begin{transitionframe}
\begin{center}
  \begin{huge}
    \textcolor{NTHU1}{%
Berrebi, Klor (2008): \\ Are Voters Sensitive to Terrorism? Direct Evidence from Israeli Electorate [student presentation]
    }
  \end{huge}
\end{center}
\end{transitionframe}

\begin{frame}
\frametitle{Takeaways and remaining questions} 
Takeaways
\begin{itemize}\setlength\itemsep{3pt}
\item Conflicts can hurt economic and institutional development
\begin{itemize}\setlength\itemsep{3pt}
\item It could shift to illicit sectors (i.e. increase undetectable activities)
  \item But the volume of economic activities can decrease
  \item State apparatus captured to bend to will of the insurgents/organized criminals
\end{itemize}
\item However, these can be reversed with right actions
\begin{itemize}\setlength\itemsep{3pt}
\item Shake-up institutional factors that cause conflicts (intergroup distrust and incentives)
\item Reduced costs of protest can lead to institutional progress
\end{itemize}
\end{itemize}\medskip
Remaining questions
\begin{itemize}\setlength\itemsep{3pt}
\item How do individuals respond to or mitigate these effects?
\begin{itemize}\setlength\itemsep{3pt}
\item Some choose to protest
\item Others migrate to seek new opportunities
\end{itemize}
\item Which types of institutions are vulnerable or resilient to these damages?
\begin{itemize}\setlength\itemsep{3pt}
\item Need to decompose state capacity and characterize ``weak states"
\end{itemize}
\end{itemize}
\end{frame}


%%%%%%%%%%%
\end{document}


